%====== FONTE E ESPAÇAMENTO ======%
\setmainfont[Contextuals={WordInitial,WordFinal}, Scale=1]{Arial}
\setmathfont{Latin Modern Math}

\onehalfspacing
\setlength{\parindent}{1.25cm}
\emergencystretch 3em

\pagestyle{fancy}
\fancyhf{} % limpa cabeçalhos e rodapés

\fancyhead[RO]{\fontsize{10pt}{12pt}\selectfont\thepage} % Right Odd
\fancyhead[LE]{\fontsize{10pt}{12pt}\selectfont\thepage} % Left Even

\renewcommand{\headrulewidth}{0pt}  % sem linha no cabeçalho
\renewcommand{\footrulewidth}{0pt}  % sem linha no rodapé

%====== SUMÁRIO ABNT ======%
% (se quiser reativar o estilo das seções principais, descomente a linha abaixo)
% \renewcommand{\cftsecfont}{\bfseries\uppercase\fontsize{12}{15}\selectfont}
\renewcommand{\cftsubsecfont}{\bfseries\fontsize{12}{15}\selectfont}
\renewcommand{\cftsubsubsecfont}{\bfseries\itshape\fontsize{12}{15}\selectfont}
\renewcommand{\cftparafont}{\itshape\fontsize{12}{15}\selectfont}
\renewcommand{\cftsubparafont}{\normalfont\fontsize{12}{15}\selectfont}
\setlength{\cftsecnumwidth}{4.5em}
\setlength{\cftsubsecnumwidth}{4.5em}
\setlength{\cftsubsubsecnumwidth}{4.5em}
\setlength{\cftparanumwidth}{4.5em}
\setlength{\cftsubparanumwidth}{4.5em}
\setlength{\cftsecindent}{0em}
\setlength{\cftsubsecindent}{0em}
\setlength{\cftsubsubsecindent}{0em}
\setlength{\cftparaindent}{0em}
\setlength{\cftsubparaindent}{0em}
\setcounter{tocdepth}{5}
\setcounter{secnumdepth}{5}

%====== TÍTULO DAS REFERÊNCIAS ======%
\renewcommand{\refname}{\MakeUppercase{referências}}
\appto{\bibsetup}{\justifying}

%====== DESTAQUES ======%
\sethlcolor{yellow}

%========    FUNÇÃO SENO    =========%
% Substituir "sin" por "sen"
\renewcommand{\sin}{\operatorname{sen}}

%====== LISTAS DE FIGURAS, QUADROS E GRÁFICOS ======%
\DeclareCaptionLabelSeparator{traco}{ - }
% Definir separadores para cada tipo
\captionsetup[figure]{labelformat=default, labelsep=traco, name=Figura}
\captionsetup[table]{labelformat=default, labelsep=traco, name=Tabela}     
            
\newlistof{grafico}{logr}{Lista de Gráficos}

\DeclareFloatingEnvironment[
  fileext=logr,
  listname={},
  name=Gráfico
]{grafico}

\captionsetup[grafico]{labelformat=default, labelsep=traco, name=Gráfico}


\makeatletter 
\renewcommand{\cftfigpresnum}{Figura~}
\renewcommand{\cftfigaftersnum}{~-~}
\renewcommand{\cftfignumwidth}{2cm}
\setlength{\cftfigindent}{0em}

\renewcommand{\listtablename}{}                        % esconde o nome na saída
\renewcommand{\cfttabpresnum}{Tabela~}                 % prefixo no sumário interno
\renewcommand{\cfttabaftersnum}{~-~}                   % sufixo “ – ” 
\setlength{\cfttabnumwidth}{2cm}     
\setlength{\cfttabindent}{0em}

\renewcommand{\listofgrafico}{%
  \begingroup
    \renewcommand{\cftloftitlefont}{\relax}
    \renewcommand{\cftafterloftitle}{\hfill}
    \renewcommand{\cftgraficopresnum}{Gráfico~}
    \renewcommand{\cftgraficoaftersnum}{~-~}
    \renewcommand{\cftgraficonumwidth}{2.2cm}
    \@starttoc{logr}
  \endgroup
}

\@ifundefined{l@grafico}{
  \newcommand*\l@grafico{\@dottedtocline{1}{1.5em}{3.2em}}
}{}

\@ifundefined{cftgraficonumwidth}{
  \newlength{\cftgraficonumwidth}
}{}

% \providecommand{\cftgraficopresnum}{Gráfico~}
% \providecommand{\cftgraficoaftersnum}{~-~}
% \setlength{\cftgraficonumwidth}{4cm}

% \renewcommand{\cftgraficoaftersnum}{\hspace{3cm}} % opcional para dar mais espaço visual

% % Alinha o número à direita dentro do espaço reservado
% \newcommand{\cftgraficoformatnum}[1]{\makebox[\cftgraficonumwidth][r]{#1}}


\makeatother